\thispagestyle{empty} 
\chapter*{Abstract}

This report presents the development of a Bell Inequality hardware accelerator and a Python-based Graphical User Interface (GUI) using the qudi framework. The internship was conducted in two parts. The first part involved implementing a hardware accelerator for verifying Bell’s Inequality, primarily through \acs{FPGA}-based VHDL design and simulation. The hardware module utilized finite state machines (FSM) to coordinate complex computation tasks essential for the experiment, ultimately enabling high-speed and reliable data processing. Additionally, the accelerator integrates a custom Division entity, designed to handle the division of experimental data in real-time with precision and efficiency. Through this process, key insights were gained into \acs{FPGA} architectures, VHDL programming, and embedded system optimizations.

The second part focused on developing a Python-based \acs{GUI} leveraging the open-source qudi platform, providing users the flexibility to implement custom Python programs and control devices such as \acs{quEDU} and \acs{quNV}. This phase introduced complex Python frameworks, modular software design, and concepts related to nitrogen-vacancy centers in diamond, which are used for quantum experiments. The software architecture was built to interface seamlessly between the \acs{FPGA} core and the \acs{GUI} layer, providing a user-friendly environment for educational and research-based quantum applications. The report further explores the interdisciplinary integration of quantum mechanics, embedded systems, and software engineering within a real-world industrial setting, highlighting the technical and theoretical learning outcomes.

In this internship, I worked extensively through the five key project phases: requirements engineering, design, implementation, testing, and delivery. While developing the Python \acs{GUI}, I gained a comprehensive understanding of requirement engineering, design, implementation and testing, balancing usability with functional requirements. In contrast, the Bell Inequality hardware accelerator provided significant experience with design and implementation, particularly through hands-on \acs{FPGA} programming and system optimization. Mastery of each phase across both projects contributed substantially to personal and professional growth, and the exposure to quantum physics and quantum system applications sparked a strong interest in advanced scientific and engineering concepts that will guide my future endeavors.

